%% kapa: exam=1 round=24 year=2013 period=1 form=A subject=civil_law subject_ko=민법 seq=26 number=26
%% source: Q-Net 감정평가사 기출문제 (cst00309) artlSeq=5073802 file=extracted/1008/5073802/2013년 제24회 감정평가사 1차시험 기출문제/1교시(A형).hwp
%% conversion: kordoc 4.24.51 -> tools/gen1.py (auto); review=proofread-sonnet; flags=none
\begin{question}{26}
부동산의 소유권취득에 관한 설명으로 옳지 않은 것은? (다툼이 있으면 판례에 의함)
\begin{choices}
  \item 자기의 비용과 노력으로 건물을 신축한 자는 원칙적으로 그 건물의 소유권을 원시취득한다.
  \item 건물신축공사의 수급인이 자기의 노력과 재료를 들여 건물을 완성하더라도 건물의 소유권을 도급인에게 귀속시키기로 하는 합의가 있는 경우 도급인이 그 건물의 소유권을 원시취득한다.
  \item 채무의 담보를 위하여 채무자가 자기 비용과 노력으로 신축하는 건물의 건축허가 명의를 채권자 명의로 하기로 합의한 경우 그 건물의 소유권은 채권자가 원시취득한다.
  \item 건축공사가 중단되어 건물의 요건을 갖추지 못한 미완성의 건물을 인도받아 자기의 비용과 노력으로 완공한 수급인은 특별한 사정이 없으면 그 건물을 원시취득한다.
  \item 매매를 원인으로 소유권이전등기절차를 이행하라는 판결이 있더라도 그 등기가 된 때에 비로소 소유권이전의 효력이 생긴다.
\end{choices}

eviewflag{question 26 stem/choices overlap in rhwp render; only last words verified}
\end{question}
\answer{3}
